\section{Design}
\label{sec:design}

\sys is placed between the agent runtime and the LLM engine (\Cref{fig:system}).
The harness announces structure; the control plane admits, prunes, binds, and retains; the backend moves pages and tokens.

\subsection{API}
\label{sec:overview}
\begin{center}
\footnotesize
\setlength{\tabcolsep}{3pt}
\begin{tabular}{@{}lp{0.56\columnwidth}@{}}
\toprule
\texttt{open(session)} & Create a root; prefill system+user. \\
\texttt{fork(u, bid, tok?)} & Child of $u$; optional known suffix. \\
\texttt{speculate(v, pri)} & Mark $v$ eligible for speculative prefill. \\
\texttt{commit(u, tokens)} & Bind by LCP; start decode. \\
\texttt{join(children, pol)} & Merge fan-out (\texttt{all}/\texttt{first}/\texttt{kofn}). \\
\texttt{abort(v)} / \texttt{close} & Drop a subtree; end the session. \\
\bottomrule
\end{tabular}
\end{center}
\texttt{commit} is the only verb on the critical path of user-visible tokens.

\subsection{Forkable Branch State}
\label{sec:state}

\begin{definition}[Context tree]
A session is a rooted tree $T=(V,E)$.
Node $u$ stores a token sequence $x_u$ such that $x_{\pi(u)}\preceq x_u$, a residual $\delta_u=x_u[|x_{\pi(u)}|{:}]$, a mode $\rho_u\in\{\textsc{Spec},\textsc{Commit},\textsc{Idle},\textsc{Dead}\}$, a generation $g_u$, and a logical page table
$\sigma_u=(\mathrm{parent},\;\mathrm{alias\_len},\;\rho^{\mathrm{pages}}_u)$.
\end{definition}

Invariants: (i)~prefix; (ii)~at most one committed child per node except during join; (iii)~speculative nodes are invisible to sampling; (iv)~a page's refcount equals the number of non-dead nodes that map it.
\texttt{fork}$(u,\mathit{bid},x^{\mathrm{k}})$ copies two words $(\mathrm{parent}\leftarrow\sigma_u,\;\mathrm{alias\_len}\leftarrow|x_u|)$ and increfs the aliased frames---$O(1)$ in $L$.
A write that hits $\mathrm{ref}>1$ or a read-only tail CoW-copies only the valid rows.
\Cref{fig:pool} is that pool.
A multi-level fan-out is a \emph{sequence} of CoWs (\Cref{fig:tree}): each edge aliases the parent and allocates only $+$r; live paths share every page on the common prefix; abort is residual-only, and a subtree with no live child cascades until the committed spine.
The two compose: LCP commit publishes full winner pages into the APC index so a later session can hash-hit them (\S\ref{sec:app}).
Abort frees $b\ell_i$ bytes, not $b(L+\ell_i)$.
TTL and offload are node-granular: speculative leaves first, then idle committed leaves, then the spine~\cite{continuum,mori}.

\begin{figure*}[t]
\centering
\resizebox{\textwidth}{!}{% Context-tree page pool. Logical verbs left, physical frames right.
\tikzset{
  fmbadge/.style={circle,draw=#1!70!black,fill=#1,text=white,
    font=\tiny\sffamily\bfseries,inner sep=0pt,minimum size=0.32cm},
  fmcard/.style={draw=#1!65,fill=#1!8,rounded corners=1.6pt,align=center,
    inner sep=2pt,font=\tiny\sffamily,line width=0.45pt},
  pg/.style 2 args={draw=#1,fill=#2,rounded corners=1.3pt,align=center,
    font=\tiny\sffamily,inner sep=1pt,minimum width=1.05cm,
    minimum height=0.70cm,line width=0.4pt},
  fp/.style={draw=fsgray,fill=black!5,rounded corners=1.2pt,align=center,
    font=\tiny\sffamily,inner sep=1pt,minimum width=0.72cm,
    minimum height=0.38cm,line width=0.35pt,text=black!70},
  note/.style={font=\tiny\sffamily,text=black!58},
  hdr/.style={font=\scriptsize\sffamily\bfseries},
  fmdata/.style={-{Latex[length=1.4mm]},line width=0.5pt,draw=fsblue!75},
  fmfeed/.style={densely dashed,-{Latex[length=1.4mm]},line width=0.5pt},
  flink/.style={{Latex[length=1.1mm,width=0.8mm]}-{Latex[length=1.1mm,width=0.8mm]},
    line width=0.36pt,draw=fsgray},
  tbox/.style={draw=#1!70!black,fill=#1!10,rounded corners=1.2pt,align=center,
    font=\tiny\sffamily,inner sep=1.2pt,line width=0.4pt},
}
\begin{tikzpicture}[x=1cm,y=1cm,font=\sffamily,
  >=Latex,line cap=round,line join=round]
\path (-0.02,-0.32) rectangle (14.72,5.78);
% ========================== zones ==========================
\fill[black!3.5,rounded corners=5pt] (0.04,0.08) rectangle (6.08,5.52);
\fill[fsblue!5,rounded corners=5pt] (6.24,0.08) rectangle (14.58,5.52);
\draw[black!22,dashed,dash pattern=on 1.2pt off 1.4pt,line width=0.5pt]
  (6.16,0.18) -- (6.16,5.38);
\node[anchor=west,hdr,text=fsgray] at (0.16,5.38) {LOGICAL \(\cdot\) TREE VERBS};
\node[anchor=west,hdr,text=fsblue] at (6.40,5.38) {PHYSICAL \(\cdot\) FRAME POOL};
% ===================== 1. request table =====================
\fill[white,rounded corners=2.5pt] (0.18,4.18) rectangle (5.92,5.28);
\fill[fsblue!16,rounded corners=2.5pt] (0.18,4.96) rectangle (5.92,5.28);
\fill[fsblue!16] (0.18,4.96) rectangle (5.92,5.08);
\draw[fsblue!65,rounded corners=2.5pt,line width=0.5pt] (0.18,4.18) rectangle (5.92,5.28);
\node[fmbadge=fsblue] at (0.42,5.12) {1};
\node[anchor=west,font=\tiny\sffamily\bfseries] at (0.64,5.12) {Request table \(\sigma_u\)};
\node[tbox=fsblue,minimum width=0.64cm,minimum height=0.30cm] (uu) at (1.22,4.74) {\(u\)};
\node[tbox=fsorange,minimum width=0.70cm,minimum height=0.30cm] (v1) at (0.52,4.40) {\(\sigma_1\)};
\node[tbox=fsgreen,minimum width=0.70cm,minimum height=0.30cm] (vs) at (1.36,4.40) {\(\sigma_\star\)};
\node[tbox=fsgray,minimum width=0.56cm,minimum height=0.30cm] (vk) at (2.12,4.40) {\(\cdots\)};
\draw[fsblue!70,densely dashed,line width=0.45pt] (uu.south) -- (v1.north);
\draw[fsblue!70,densely dashed,line width=0.45pt] (uu.south) -- (vs.north);
\draw[fsblue!70,densely dashed,line width=0.45pt] (uu.south) -- (vk.north);
\node[fmcard=fsblue,anchor=west,align=left,text width=2.85cm,minimum height=0.62cm]
  at (2.62,4.56)
  {\(\sigma_u{=}(\mathrm{parent},\,\mathrm{alias\_len},\,\rho)\)\\[1pt]
   \(\mathrm{ref}(p){=}\#\) live mappers};
% ===================== 2. fork =====================
\fill[white,rounded corners=2.5pt] (0.18,3.02) rectangle (5.92,4.04);
\fill[fsblue!16,rounded corners=2.5pt] (0.18,3.72) rectangle (5.92,4.04);
\fill[fsblue!16] (0.18,3.72) rectangle (5.92,3.84);
\draw[fsblue!65,rounded corners=2.5pt,line width=0.5pt] (0.18,3.02) rectangle (5.92,4.04);
\node[fmbadge=fsblue] at (0.42,3.88) {2};
\node[anchor=west,font=\tiny\sffamily\bfseries] at (0.64,3.88)
  {\texttt{fork}\((u)\): two words, \(O(1)\) in \(L\)};
\node[fmcard=fsblue,minimum width=1.55cm,minimum height=0.40cm] (cp) at (1.15,3.42)
  {\(\mathrm{parent}\leftarrow\sigma_u\)};
\node[fmcard=fsblue,minimum width=1.62cm,minimum height=0.40cm] (ca) at (2.95,3.42)
  {\(\mathrm{alias\_len}\leftarrow|x_u|\)};
\node[fmcard=fsorange,minimum width=1.35cm,minimum height=0.40cm] (ci) at (4.85,3.42)
  {incref trunk};
\draw[fmdata] (cp) -- (ca);
\draw[fmdata] (ca) -- (ci);
\node[note] at (3.05,3.16) {no memcpy of \(L\)};
% ===================== 6. LCP =====================
\fill[white,rounded corners=2.5pt] (0.18,1.52) rectangle (5.92,2.88);
\fill[fsgreen!18,rounded corners=2.5pt] (0.18,2.56) rectangle (5.92,2.88);
\fill[fsgreen!18] (0.18,2.56) rectangle (5.92,2.68);
\draw[fsgreen!60!black,rounded corners=2.5pt,line width=0.5pt] (0.18,1.52) rectangle (5.92,2.88);
\node[fmbadge=fsgreen] at (0.42,2.72) {6};
\node[anchor=west,font=\tiny\sffamily\bfseries] at (0.64,2.72) {LCP commit};
\node[fmcard=fsgreen,minimum width=5.35cm,minimum height=0.34cm] at (3.05,2.36)
  {\(v^\star{=}\arg\max_v\mathrm{LCP}(x,x_v)\)};
\node[fmcard=fsgreen,minimum width=5.35cm,minimum height=0.34cm] at (3.05,1.98)
  {\(H\leftarrow H\cup\{x_\star\}\)};
\node[fmcard=fspurple,minimum width=5.35cm,minimum height=0.34cm] at (3.05,1.66)
  {\(g\leftarrow g{+}1\) retires in-flight spec};
% ===================== 7. abort =====================
\fill[white,rounded corners=2.5pt] (0.18,0.16) rectangle (5.92,1.38);
\fill[fsorange!18,rounded corners=2.5pt] (0.18,1.06) rectangle (5.92,1.38);
\fill[fsorange!18] (0.18,1.06) rectangle (5.92,1.18);
\draw[fsorange!70,rounded corners=2.5pt,line width=0.5pt] (0.18,0.16) rectangle (5.92,1.38);
\node[fmbadge=fsorange] at (0.42,1.22) {7};
\node[anchor=west,font=\tiny\sffamily\bfseries] at (0.64,1.22) {\texttt{abort}};
\node[fmcard=fsorange,minimum width=5.35cm,minimum height=0.46cm] at (3.05,0.62)
  {\(M_{\mathrm{spine}}{=}b(L{+}\ell_\star)\)\\[1pt] free \(b\ell_i\), not \(b(L{+}\ell_i)\)};
% ===================== 3. frames =====================
\fill[fsblue!7,rounded corners=3pt] (6.38,3.02) rectangle (14.44,5.28);
\draw[fsblue!45,rounded corners=3pt,line width=0.5pt] (6.38,3.02) rectangle (14.44,5.28);
\node[fmbadge=fsblue] at (6.60,5.08) {3};
\node[anchor=west,font=\tiny\sffamily] at (6.82,5.08) {\textbf{Shared trunk, private \(\rho\)}};
\node[pg={fsblue}{fslight},minimum width=2.35cm,minimum height=0.92cm] at (7.75,4.22)
  {\textbf{shared trunk}\\[-1pt] one copy of \(L\)\\[-1pt] read-only};
\node[pg={fsorange!80!black}{fsspec},minimum width=2.35cm,minimum height=0.92cm] at (10.35,4.22)
  {\textbf{private \(\rho\)}\\[-1pt] new frames only\\[-1pt] \(\mathrm{ref}{=}1\), not in \(H\)};
\node[pg={fsorange!80!black}{fsspec},densely dashed,minimum width=2.35cm,minimum height=0.92cm] at (12.95,4.22)
  {\textbf{partial tail}\\[-1pt] last page \(n{<}P\)\\[-1pt] not published};
\node[note] at (10.41,3.42)
  {child aliases \(L\);\ a write with \(\mathrm{ref}{>}1\) copies that page's valid rows};
% ===================== 4. hash index =====================
\fill[fsgreen!8,rounded corners=3pt] (6.38,1.52) rectangle (14.44,2.88);
\draw[fsgreen!50!black,rounded corners=3pt,line width=0.5pt] (6.38,1.52) rectangle (14.44,2.88);
\node[fmbadge=fsgreen] at (6.60,2.68) {4};
\node[anchor=west,font=\tiny\sffamily] at (6.82,2.68) {\textbf{Publish index \(H\)}};
\node[fmcard=fsgreen,minimum width=3.55cm,minimum height=0.62cm] at (8.40,2.02)
  {in \(H\): full committed block\\[-1pt] \(h(\mathrm{parent},\mathrm{block})\)};
\node[fmcard=fsorange,minimum width=3.55cm,minimum height=0.62cm] at (12.35,2.02)
  {out of \(H\): speculative\\[-1pt] or partial page};
% ===================== 5. free list =====================
\fill[black!3,rounded corners=3pt] (6.38,0.16) rectangle (14.44,1.38);
\draw[fsgray!70,rounded corners=3pt,line width=0.5pt] (6.38,0.16) rectangle (14.44,1.38);
\node[fmbadge=fsgray] at (6.60,1.18) {5};
\node[anchor=west,font=\tiny\sffamily] at (6.82,1.18) {\textbf{Free list}\quad \(\mathrm{ref}{=}0\)};
\node[fmcard=fsblue,minimum width=3.55cm,minimum height=0.58cm] at (8.40,0.58)
  {\texttt{fork} takes free frames\\[-1pt] for the new \(\rho\)};
\node[fmcard=fsorange,minimum width=3.55cm,minimum height=0.58cm] at (12.35,0.58)
  {\texttt{abort} returns those frames\\[-1pt] trunk stays allocated};
% ===================== cross-boundary =====================
\draw[fmfeed,draw=fsblue,rounded corners=1.5pt]
  (5.92,3.42) -- (6.28,3.42) -- (6.28,4.28) -- (6.55,4.28);
\draw[fmfeed,draw=fsgreen!50!black] (5.92,1.98) -- (6.55,1.98);
\draw[fmfeed,draw=fsorange!80!black] (5.92,0.62) -- (6.55,0.62);
% ===================== legend =====================
\node[anchor=west,font=\tiny\sffamily,text=black!58] at (0.06,-0.14)
  {\textcolor{fsblue}{trunk}\,\(\cdot\)\,%
   \textcolor{fsorange}{residual}\,\(\cdot\)\,%
   \textcolor{fsgreen}{publish}\,\(\cdot\)\,%
   \textcolor{fspurple}{generation}\,\(\cdot\)\,%
   \textcolor{fsgray}{free}\qquad
   dashed edge \(=\) alias};
\end{tikzpicture}
\iffalse
% Alternate symbolic drawing. Kept for reuse; not compiled.
% Symbolic page pool. The previous drawing is kept below and is not rendered.
\begin{tikzpicture}[
  x=1cm, y=1cm,
  font=\scriptsize\sffamily,
  >=Latex,
  line cap=round, line join=round,
  panel/.style={draw=fsblue!28, line width=0.5pt, rounded corners=1.2pt, fill=white},
  head/.style={font=\scriptsize\bfseries\sffamily, text=white, anchor=west},
  eq/.style={font=\scriptsize\sffamily, anchor=west, inner sep=0.5pt},
  pg/.style={draw=fsblue!70, fill=fsblue!12, minimum width=0.72cm, minimum height=0.52cm,
             font=\tiny\sffamily, inner sep=1pt},
  priv/.style={draw=fsorange!75!black, fill=fsspec, minimum width=0.62cm, minimum height=0.42cm,
               font=\tiny\sffamily, inner sep=0.6pt},
  alias/.style={densely dashed, line width=0.4pt, draw=fsblue!60},
]
\path (0,0) rectangle (16.45,5.55);

\draw[panel] (0,0) rectangle (7.55,5.55);
\draw[panel] (7.78,0) rectangle (16.45,5.55);
\fill[fsblue] (0,5.18) rectangle (7.55,5.55);
\fill[fsorange!85!black] (7.78,5.18) rectangle (16.45,5.55);
\node[head] at (0.14,5.36) {operations};
\node[head] at (7.92,5.36) {one physical trunk, $k$ aliases};

\node[eq] at (0.18,4.78) {$\texttt{fork}(\sigma_u)$};
\node[eq] at (0.36,4.32) {$\sigma_i\leftarrow(\sigma_u,|x_u|),\quad \mathrm{ref}(p){+}{=}1$};
\node[eq] at (0.36,3.86) {allocate $\rho_i$ only,\ $O(1)$};

\node[eq] at (0.18,3.28) {write on $p$,\ $\mathrm{ref}(p)>1$};
\node[eq] at (0.36,2.82) {copy the valid rows of $p$};
\node[eq] at (0.36,2.36) {$\mathrm{ref}(p){-}{=}1$, writer owns the copy};

\node[eq] at (0.18,1.78) {$\texttt{lcp\_commit}(x_\star)$};
\node[eq] at (0.36,1.32) {$H\leftarrow H\cup\{x_\star\},\quad g\leftarrow g{+}1$};

\node[eq] at (0.18,0.74) {$\texttt{abort}(\sigma_i)$};
\node[eq] at (0.36,0.28) {free $\rho_i$;\ $\mathrm{ref}(p){=}0\Rightarrow$ free list};

% right: before abort — shared pages, private residual
\node[font=\tiny\sffamily\bfseries, anchor=west] at (7.96,4.95) {before abort};
\foreach \x in {8.55,9.40,10.25,11.10} {
  \node[pg] at (\x,4.38) {$p$};
}
\node[font=\tiny\sffamily, anchor=west, text=fsblue] at (11.62,4.38) {$\mathrm{ref}(p){=}k$};
\node[draw=fsgreen!60!black, fill=fscommit, rounded corners=1pt, font=\tiny\sffamily,
      inner sep=1.2pt, minimum width=0.95cm] (a0) at (8.70,3.28) {$\sigma_\star$};
\node[draw=fsgray, fill=black!4, densely dashed, rounded corners=1pt, font=\tiny\sffamily,
      inner sep=1.2pt, minimum width=0.95cm] (a1) at (11.15,3.28) {$\sigma_i$};
\node[draw=fsorange!75!black, fill=fsspec, rounded corners=1pt, font=\tiny\sffamily,
      inner sep=1.2pt, minimum width=0.95cm] (ak) at (13.60,3.28) {$\sigma_k$};
\draw[alias] (a0.north) -- (9.40,4.12);
\draw[alias] (a1.north) -- (9.82,4.12);
\draw[alias] (ak.north) -- (10.25,4.12);
\node[priv] (r0) at (8.70,2.52) {$\rho_\star$};
\node[priv] (r1) at (11.15,2.52) {$\rho_i$};
\node[priv] (rk) at (13.60,2.52) {$\rho_k$};
\draw[-{Latex[length=1.4mm]}, line width=0.4pt, draw=fsorange!75!black] (a0.south) -- (r0.north);
\draw[-{Latex[length=1.4mm]}, line width=0.4pt, draw=fsorange!75!black] (a1.south) -- (r1.north);
\draw[-{Latex[length=1.4mm]}, line width=0.4pt, draw=fsorange!75!black] (ak.south) -- (rk.north);
\node[font=\tiny\sffamily, text=fsgray, anchor=west] at (14.55,2.52) {$\notin H$};

% right: after
\node[font=\tiny\sffamily\bfseries, anchor=west] at (7.96,1.95) {after abort of $i\neq\star$};
\foreach \x in {8.55,9.40,10.25,11.10} {
  \node[pg] at (\x,1.38) {$p$};
}
\node[font=\tiny\sffamily, anchor=west, text=fsgreen!35!black] at (11.62,1.38) {$\mathrm{ref}(p)\ge 1$};
\node[draw=fsgreen!60!black, fill=fscommit, rounded corners=1pt, font=\tiny\sffamily,
      inner sep=1.2pt, minimum width=0.95cm] (b0) at (8.70,0.48) {$\sigma_\star$};
\draw[alias] (b0.north) -- (9.40,1.12);
\node[priv] (rb) at (10.55,0.48) {$\rho_\star$};
\draw[-{Latex[length=1.4mm]}, line width=0.4pt, draw=fsorange!75!black] (b0.east) -- (rb.west);
\node[font=\scriptsize\sffamily, anchor=west] at (11.55,0.48)
  {$M_{\mathrm{spine}}=b(L+\ell_\star)$};
\end{tikzpicture}

\fi
}
\caption{\textbf{Context-tree page pool.}
\texttt{fork} aliases the trunk in $O(1)$ and allocates only $\rho$~\textcircled{1}--\textcircled{3}.
LCP commit publishes $x_\star$ into $H$~\textcircled{4}\textcircled{6}.
\texttt{abort} returns residual frames only, so live occupancy is the spine $M_{\mathrm{spine}}{=}b(L{+}\ell_\star)$~\textcircled{5}\textcircled{7}.}
\label{fig:pool}
\end{figure*}

\subsection{Admission}
\label{sec:speculate}
Candidates $\Br(u)$ come from harness-declared wraps, grammar mass over the next tool name (top-$m$), and a session-local prior of $(\mathit{parent\_role},\mathit{bid})$.
Priors never invent a branch id.
Each candidate is a pair $(x^{\mathrm{k}}_b,\hat{x}^{\mathrm{o}}_b)$: an exact known suffix and an optional observation guess.
This is \emph{input-side} speculation---a cache of a future prompt, not draft output tokens~\cite{speculative,medusa}.

The expected \ttft{} reduction and the interference-aware cost of a work item $w$ are
\begin{align}
G(w) &= p(w)\cdot\min\bigl(\Pref(|w|),\;\Idle(u)\bigr)\cdot\eta(w), \label{eq:gain}\\
C(w) &= b\cdot|w|+\lambda\cdot\max\bigl(0,\;\Pref(|w|)-\gamma\bigr), \label{eq:cost}
\end{align}
where $\eta$ down-weights residuals by $q_b$ and work that would finish after the pause, and $\lambda$ matches $1$\,ms \tbt{} to $4$\,ms \ttft{}~\cite{jitserve}.

\begin{proposition}[Known-suffix first]
\label{prop:known}
If two items have equal $\Pref$ and $w$ is a known suffix ($p{=}1,q{=}1$) while $w'$ is an observation residual ($q<1$), then $G(w)/C(w)\ge G(w')/C(w')$.
\end{proposition}

Admission is a greedy fractional knapsack on $G/C$ subject to idle horizon, residual HBM, branch cap $m$, and committed \tbt{} margin.
Chunks of size $\le C_{\mathrm{spec}}$ make a mid-prefill \texttt{commit} discard at most one in-flight chunk.

\paragraph{LCP commit.}
When the harness presents prompt $x$ of parent $u$:
(1)~$v^\star=\arg\max_v\mathrm{LCP}(x,x_v)$ among live children;
(2)~pages covering the LCP become committed (CoW-split if mid-page);
(3)~the unmatched tail is a committed prefill;
(4)~sibling Spec nodes abort;
(5)~$g_u\leftarrow g_u+1$ and decode is released.
Speculative jobs tagged $g\neq g_u$ drop in $O(1)$.

\begin{theorem}[Output identity]
\label{thm:correct}
For any session, the committed decode tokens of \sys equal those of a baseline that prefills each committed prompt from scratch (or from a prefix cache of committed tokens only), given equal seeds and weights.
\end{theorem}
\begin{proof}[Proof sketch]
Decode at a committed node $u$ attends only to $\mathrm{KV}(x_u)$.
That KV is produced by a committed prefill of $x_u$, or by a speculative prefix verified token-for-token by LCP plus a committed tail.
Transformer KV of a prefix depends only on that prefix.
Speculative nodes do not enter the sampler.
\end{proof}

\subsection{Cost-Ordered Prefill Pruning and Transfer Gating}
\label{sec:app}

Copy-on-write fixes trunk placement.
It does not decide which residual enters a GPU kernel, nor which KV crosses the prefill--decode connector.
\app makes both decisions per residual, $(\Work_i,\Xfer_i)$, by a cascade whose stage cost is strictly increasing (\Cref{fig:app}).
The stages are a block-hash probe, a CPU draft score, a prefix of the prefill kernel, the miss-tail prefill, and a connector \texttt{insert}.
Stage $i{+}1$ runs only when stage $i$ does not reject $\rho_i$.
The committed winner is exempt: index $0$ is forced into $\Keep$, speculative pages stay unpublished until LCP commit, and Theorem~\ref{thm:correct} is unchanged.
CoW remains the storage substrate; \app is the filter on compute and on transfer.
On the Qwen3-14B GSM8K forest, copy-on-write already stores less KV than APC, yet median fan-out is $140$\,ms against APC's $87$\,ms: the $k$ residuals are still submitted to prefill.
The cascade is the stage that removes that work.

\begin{figure*}[t]
\centering
\resizebox{\textwidth}{!}{\begin{tikzpicture}[
  x=1cm, y=1cm,
  font=\sffamily,
  >=Latex,
  line cap=round, line join=round,
  box/.style={draw=fsblue!55, rounded corners=1.2pt, fill=white, inner sep=2.2pt,
              align=center, font=\tiny\sffamily, line width=0.45pt, minimum height=0.72cm},
  op/.style={circle, draw=fsblue!70, fill=fslight, inner sep=0.4pt,
             font=\tiny\sffamily\bfseries, minimum size=0.38cm, line width=0.45pt},
  lab/.style={font=\tiny\sffamily, text=fsgray, align=center},
  arr/.style={-{Latex[length=1.5mm,width=1.15mm]}, line width=0.45pt, draw=fsblue!70},
  sarr/.style={-{Latex[length=1.2mm,width=0.9mm]}, densely dashed, draw=fsgray, line width=0.35pt},
]
\path (0,-1.50) rectangle (16.50,7.15);

% ===== legend =====
\node[font=\tiny\sffamily, anchor=west, text=fsblue] at (0.05,6.95)
  {$\rho_i$: $i$-th residual\quad
   $h_i$: block hash\quad
   $s_i$: draft score\quad
   $\alpha$: early-abort fraction\quad
   $\kappa$: chunk\quad
   $\Work_i$: GPU miss tail\quad
   $\Xfer_i$: connector tokens\quad
   $a_i$: action};
\draw[fsblue!25, line width=0.3pt] (0.05,6.72) -- (16.40,6.72);

% ===== (a) APP+PD pipeline =====
\node[font=\scriptsize\bfseries\sffamily, anchor=west] at (0.05,6.48)
  {(a) \app{} + P/D pipeline (Alg.~\ref{alg:app})};

\node[box, fill=fslight, minimum width=1.55cm] (r) at (0.95,5.35)
  {$\{\rho_i\}_{i=1}^{k}$\\[-1pt]{\tiny residuals}};
\node[op] (o1) at (2.15,5.35) {$\otimes$};
\node[box, minimum width=1.85cm] (h) at (3.55,5.35)
  {Hash lookup\\[-1pt]{\tiny $h_i=\mathrm{hash}(\pi,\mathrm{blk})$}};
\node[op] (o2) at (4.85,5.35) {$\oplus$};
\node[box, minimum width=1.70cm] (d) at (6.15,5.35)
  {Draft prune\\[-1pt]{\tiny $s_i\gtrless\tau$}};
\node[op] (o3) at (7.40,5.35) {$\rhd$};
\node[box, minimum width=1.70cm] (e) at (8.70,5.35)
  {Early abort\\[-1pt]{\tiny first $\alpha\ell_i$}};
\node[op] (o4) at (9.95,5.35) {$\bullet$};
\node[box, fill=fscommit, minimum width=1.70cm] (g) at (11.30,5.35)
  {GPU miss tail\\[-1pt]{\tiny $\Work_i$}};
\node[op] (o5) at (12.55,5.35) {$\triangleright$};
\node[box, fill=fsspec, minimum width=1.55cm] (x) at (13.80,5.35)
  {\texttt{insert}\\[-1pt]{\tiny $\Xfer_i$}};
\node[box, fill=fscommit, minimum width=1.35cm] (dec) at (15.50,5.35)
  {Decode\\[-1pt]{\tiny $x_\star$}};

\draw[arr] (r) -- (o1);
\draw[arr] (o1) -- (h);
\draw[arr] (h) -- (o2);
\draw[arr] (o2) -- (d);
\draw[arr] (d) -- (o3);
\draw[arr] (o3) -- (e);
\draw[arr] (e) -- (o4);
\draw[arr] (o4) -- (g);
\draw[arr] (g) -- (o5);
\draw[arr] (o5) -- (x);
\draw[arr] (x) -- (dec);

% side exits
\node[lab, text=fsblue] at (3.55,4.55) {\texttt{HASH\_SKIP}};
\node[lab] at (6.15,4.55) {\texttt{DRAFT\_SKIP}};
\node[lab, text=fsorange!80!black] at (8.70,4.55) {\texttt{EARLY\_ABORT}};
\draw[sarr] (h.south) -- ++(0,-0.42);
\draw[sarr] (d.south) -- ++(0,-0.42);
\draw[sarr] (e.south) -- ++(0,-0.42);

% annotations under (a)
\node[lab, anchor=west] at (0.10,4.05)
  {Parallelism: $k$ residuals \quad
   Layout: location-addressed CoW, not content-addressed APC \quad
   Winner $i{=}0$ is never dropped \quad
   Spec pages unpublished};
\node[lab, anchor=west] at (0.10,3.72)
  {Work $\Work{=}\sum_i\Work_i$ \quad
   Transfer $\Xfer{=}\sum_{i\in\Keep}\Work_i$ \quad
   Occupancy after abort $M_{\mathrm{spine}}=b(L+\ell_\star)$ \quad
   Replay of $x_\star$: $\Work{=}\Xfer{=}0$};

% ===== (b) four dataflows =====
\node[font=\scriptsize\bfseries\sffamily, anchor=west] at (0.05,3.32)
  {(b) First-fan-out dataflow (same $k,L,\ell$). Bottleneck in red.};

% row helper coords
% APC
\node[font=\tiny\sffamily\bfseries, anchor=east, text=fsblue] at (1.55,2.72) {APC};
\node[box, minimum width=1.35cm, minimum height=0.55cm] (a1) at (2.45,2.72) {$k$ reqs};
\node[box, minimum width=1.55cm, minimum height=0.55cm, fill=fsmiss] (a2) at (4.35,2.72) {hash miss};
\node[box, minimum width=1.85cm, minimum height=0.55cm, fill=fsmiss] (a3) at (6.55,2.72) {clone $k$ trunks};
\node[box, minimum width=1.55cm, minimum height=0.55cm] (a4) at (8.65,2.72) {decode};
\draw[arr] (a1) -- (a2); \draw[arr] (a2) -- (a3); \draw[arr] (a3) -- (a4);
\node[lab, anchor=west] at (9.60,2.72)
  {$\Work{=}kL{+}\sum\ell_i$ \quad $\Xfer{=}0$ \quad $M{=}M_{\mathrm{clone}}$};

% stock PD
\node[font=\tiny\sffamily\bfseries, anchor=east, text=fspurple] at (1.55,2.00) {P/D disagg.};
\node[box, minimum width=1.35cm, minimum height=0.55cm] (p1) at (2.45,2.00) {$k$ reqs};
\node[box, minimum width=1.55cm, minimum height=0.55cm, fill=fsmiss] (p2) at (4.35,2.00) {prefill all};
\node[box, minimum width=1.85cm, minimum height=0.55cm, fill=fsmiss] (p3) at (6.55,2.00) {\texttt{insert} all};
\node[box, minimum width=1.55cm, minimum height=0.55cm] (p4) at (8.65,2.00) {\texttt{drop\_select}};
\draw[arr] (p1) -- (p2); \draw[arr] (p2) -- (p3); \draw[arr] (p3) -- (p4);
\node[lab, anchor=west] at (9.60,2.00)
  {$\Work{=}kL{+}\sum\ell_i$ \quad $\Xfer{=}kL{+}\sum\ell_i$ \quad $M{=}M_{\mathrm{clone}}$};

% ForkServe CoW
\node[font=\tiny\sffamily\bfseries, anchor=east, text=fsorange!80!black] at (1.55,1.28) {FS CoW};
\node[box, minimum width=1.35cm, minimum height=0.55cm, fill=fslight] (f1) at (2.45,1.28) {\texttt{fork}};
\node[box, minimum width=1.55cm, minimum height=0.55cm] (f2) at (4.35,1.28) {alias $L$};
\node[box, minimum width=1.85cm, minimum height=0.55cm, fill=fsspec] (f3) at (6.55,1.28) {prefill $k$ $\rho_i$};
\node[box, minimum width=1.55cm, minimum height=0.55cm] (f4) at (8.65,1.28) {abort};
\draw[arr] (f1) -- (f2); \draw[arr] (f2) -- (f3); \draw[arr] (f3) -- (f4);
\node[lab, anchor=west] at (9.60,1.28)
  {$\Work{=}L{+}\sum\ell_i$ \quad $\Xfer{=}0$ \quad $M{\to}M_{\mathrm{spine}}$};

% APP+PD
\node[font=\tiny\sffamily\bfseries, anchor=east, text=fsgreen!40!black] at (1.55,0.56) {\app{}+P/D};
\node[box, minimum width=1.35cm, minimum height=0.55cm, fill=fslight] (n1) at (2.45,0.56) {\texttt{fork}};
\node[box, minimum width=1.55cm, minimum height=0.55cm, fill=fscommit] (n2) at (4.35,0.56) {\app{} filter};
\node[box, minimum width=1.85cm, minimum height=0.55cm, fill=fscommit] (n3) at (6.55,0.56) {\texttt{insert} $\Keep$};
\node[box, minimum width=1.55cm, minimum height=0.55cm, fill=fscommit] (n4) at (8.65,0.56) {spine};
\draw[arr] (n1) -- (n2); \draw[arr] (n2) -- (n3); \draw[arr] (n3) -- (n4);
\node[lab, anchor=west] at (9.60,0.56)
  {$\Work{=}L{+}\sum_{\Keep}\ell_i$ \quad $\Xfer{=}\sum_{\Keep}\Work_i$ \quad $M{\to}M_{\mathrm{spine}}$};

\node[font=\tiny\sffamily\bfseries, anchor=east, text=fsgreen!40!black] at (1.55,-0.16) {overlap};
\node[box, minimum width=1.35cm, minimum height=0.55cm, fill=fslight] (v1) at (2.45,-0.16) {keep $\rho_\star$};
\node[box, minimum width=1.55cm, minimum height=0.55cm, fill=fscommit] (v2) at (4.35,-0.16) {decode step};
\node[box, minimum width=1.85cm, minimum height=0.55cm, fill=fscommit] (v3) at (6.55,-0.16) {prefill next};
\node[box, minimum width=1.55cm, minimum height=0.55cm] (v4) at (8.65,-0.16) {admit slot};
\draw[arr] (v1) -- (v2); \draw[arr] (v2) -- (v3); \draw[arr] (v3) -- (v4);
\node[lab, anchor=west] at (9.60,-0.16)
  {$B^s_t$ prefills the next winner \quad losers are not in the queue};

\node[font=\tiny\sffamily\bfseries, anchor=east, text=fsorange!80!black] at (1.55,-0.88) {tool idle};
\node[box, minimum width=1.35cm, minimum height=0.55cm, fill=fsspec] (i1) at (2.45,-0.88) {pause};
\node[box, minimum width=1.55cm, minimum height=0.55cm, fill=fsspec] (i2) at (4.35,-0.88) {pin $p{=}1$};
\node[box, minimum width=1.85cm, minimum height=0.55cm, fill=fscommit] (i3) at (6.55,-0.88) {obs arrives};
\node[box, minimum width=1.55cm, minimum height=0.55cm] (i4) at (8.65,-0.88) {tail only};
\draw[arr] (i1) -- (i2); \draw[arr] (i2) -- (i3); \draw[arr] (i3) -- (i4);
\node[lab, anchor=west] at (9.60,-0.88)
  {$\Pref\le\Idle(u)$ \quad else reject \quad $C{=}\lfloor H/M_{\mathrm{spine}}\rfloor$};

\node[lab, anchor=west] at (0.10,-1.28)
  {Takeaway: pruning removes non-winner work; the spine then sets the slot count, shares the decode step with the next winner, and absorbs a known suffix during the tool pause.};
\end{tikzpicture}
}
\caption{\textbf{Cost-ordered prefill pruning and the critical path.}
(a)~A residual exits at the first decisive stage. Only a kept miss tail reaches \texttt{insert}.
(b)~APC and baseline disaggregation clone the trunk. CoW aliases $L$. \app{} keeps the winner.
The overlap row prefills that next winner in the decode step. The tool-idle row pins a known suffix during the pause and leaves the observation tail.}
\label{fig:app}
\end{figure*}

\begin{figure*}[t]
\centering
\resizebox{\textwidth}{!}{\begin{tikzpicture}[
  font=\scriptsize,
  >=Stealth,
  live/.style={draw=fsgreen!70!black, fill=fscommit, rounded corners=1.3pt,
               inner sep=1.3pt, align=center, font=\tiny\sffamily, minimum height=0.38cm},
  dead/.style={draw=fsgray, fill=black!7, rounded corners=1.3pt, densely dashed,
               inner sep=1.3pt, align=center, font=\tiny\sffamily, minimum height=0.38cm, text=fsgray},
  casc/.style={draw=fspurple, fill=fspurple!12, rounded corners=1.3pt, densely dashed,
               inner sep=1.3pt, align=center, font=\tiny\sffamily, minimum height=0.38cm, text=fspurple},
  root/.style={draw=fsblue, fill=fslight, rounded corners=1.5pt,
               inner sep=1.6pt, align=center, font=\tiny\sffamily\bfseries, minimum height=0.44cm},
  arr/.style={-{Stealth[length=2.4pt,width=1.8pt]}, thick, fsblue!80},
  darr/.style={-{Stealth[length=2.4pt,width=1.8pt]}, thick, densely dashed, fsorange!80},
  carr/.style={-{Stealth[length=2.4pt,width=1.8pt]}, thick, densely dashed, fspurple},
  rlab/.style={font=\tiny\ttfamily, text=fsorange!80!black},
  sc/.style={font=\tiny, text=fsgreen!40!black},
  sd/.style={font=\tiny, text=fsorange!75!black},
  stg/.style={font=\tiny\sffamily, text=fsblue!70, align=right, anchor=east},
]
% ============================================================
% (a)
% ============================================================
\node[font=\small\bfseries, anchor=west] at (-0.35,6.28) {(a) Staged CoW: one residual per fork, shared prefix across paths};
\node[font=\tiny, text=fsgray, anchor=west] at (-0.35,5.92)
  {abort iff $s{<}\tau$ ($\tau{=}0.40$). Fork C (depth 2): leaves die in one round $\Rightarrow$ parents with no live child cascade up.};

% stage labels
\node[stg] at (-1.72,5.38) {trunk};
\node[stg] at (-1.72,4.32) {CoW-1};
\node[stg] at (-1.72,3.28) {CoW-2};
\node[stg] at (-1.72,2.24) {CoW-3};
\node[stg] at (-1.72,1.20) {CoW-4};
\node[stg] at (-1.72,0.08) {CoW-5};

\node[root] (R) at (4.85,5.38) {Root / Planner \\[-1pt]{\tiny one shared trunk}};

% CoW-1: A, B, and cascade root C
\node[live] (A) at (1.55,4.32) {A};
\node[live] (B) at (4.85,4.32) {B};
\node[casc] (C) at (9.15,4.32) {C $\Uparrow$};
\node[rlab] at (1.55,3.94) {$+$r$_A$};
\node[rlab] at (4.85,3.94) {$+$r$_B$};
\node[rlab] at (9.15,3.94) {$+$r$_C$};
\node[sc] at (2.22,4.32) {$s{=}0.82$};
\node[sc] at (5.52,4.32) {$s{=}0.71$};

% CoW-2
\node[live] (A1) at (0.40,3.28) {A1};
\node[dead] (A2) at (1.55,3.28) {A2 $\times$};
\node[live] (A3) at (2.70,3.28) {A3};
\node[live] (B2) at (4.40,3.28) {B2};
\node[dead] (B1) at (5.55,3.28) {B1 $\times$};
\node[casc] (C1) at (8.25,3.28) {C1 $\Uparrow$};
\node[casc] (C2) at (10.05,3.28) {C2 $\Uparrow$};
\node[rlab] at (0.40,2.90) {$+$r$_{A1}$};
\node[rlab] at (2.70,2.90) {$+$r$_{A3}$};
\node[rlab] at (4.40,2.90) {$+$r$_{B2}$};
\node[sd] at (1.55,2.90) {$0.31{<}\tau$};
\node[sd] at (5.55,2.90) {$0.22{<}\tau$};
\node[rlab] at (8.25,2.90) {$+$r$_{C1}$};
\node[rlab] at (10.05,2.90) {$+$r$_{C2}$};

% CoW-3: original plus C's depth-2 leaves (all die this round)
\node[dead] (A1a) at (-0.70,2.24) {A1a $\times$};
\node[live] (A1b) at (0.40,2.24) {A1b};
\node[live] (A3a) at (2.70,2.24) {A3a};
\node[live] (B2a) at (4.40,2.24) {B2a};
\node[dead] (C1a) at (7.55,2.24) {C1a $\times$};
\node[dead] (C1b) at (8.70,2.24) {C1b $\times$};
\node[dead] (C2a) at (9.55,2.24) {C2a $\times$};
\node[dead] (C2b) at (10.70,2.24) {C2b $\times$};
\node[rlab] at (0.40,1.86) {$+$r$_{A1b}$};
\node[rlab] at (2.70,1.86) {$+$r$_{A3a}$};
\node[rlab] at (4.40,1.86) {$+$r$_{B2a}$};
\node[sd] at (-0.70,1.86) {$0.35{<}\tau$};
\node[sd] at (7.55,1.86) {$0.24$};
\node[sd] at (8.70,1.86) {$0.28$};
\node[sd] at (9.55,1.86) {$0.17$};
\node[sd] at (10.70,1.86) {$0.22$};

% CoW-4
\node[live] (A1b1) at (0.40,1.20) {A1b1};
\node[live] (A3a2) at (2.70,1.20) {A3a2};
\node[live] (B2a1) at (4.40,1.20) {B2a1};
\node[rlab] at (0.40,0.82) {$+$r$_{A1b1}$};
\node[rlab] at (2.70,0.82) {$+$r$_{A3a2}$};
\node[rlab] at (4.40,0.82) {$+$r$_{B2a1}$};

% CoW-5
\node[live] (A1b1x) at (-0.25,0.08) {A1b1x};
\node[dead] (A1b1y) at (1.05,0.08) {A1b1y $\times$};
\node[live] (A3a2x) at (2.70,0.08) {A3a2x};
\node[dead] (B2a1x) at (3.90,0.08) {B2a1x $\times$};
\node[live] (B2a1y) at (5.10,0.08) {B2a1y};
\node[rlab] at (-0.25,-0.30) {$+$r};
\node[sd] at (1.05,-0.30) {$0.27{<}\tau$};
\node[rlab] at (2.70,-0.30) {$+$r};
\node[sd] at (3.90,-0.30) {$0.33{<}\tau$};
\node[rlab] at (5.10,-0.30) {$+$r};

\draw[arr] (R) -- (A);
\draw[arr] (R) -- (B);
\draw[carr] (R) -- (C);
\draw[arr] (A) -- (A1);
\draw[darr] (A) -- (A2);
\draw[arr] (A) -- (A3);
\draw[arr] (B) -- (B2);
\draw[darr] (B) -- (B1);
\draw[carr] (C) -- (C1);
\draw[carr] (C) -- (C2);
\draw[darr] (A1) -- (A1a);
\draw[arr] (A1) -- (A1b);
\draw[arr] (A3) -- (A3a);
\draw[arr] (B2) -- (B2a);
\draw[darr] (C1) -- (C1a);
\draw[darr] (C1) -- (C1b);
\draw[darr] (C2) -- (C2a);
\draw[darr] (C2) -- (C2b);
\draw[arr] (A1b) -- (A1b1);
\draw[arr] (A3a) -- (A3a2);
\draw[arr] (B2a) -- (B2a1);
\draw[arr] (A1b1) -- (A1b1x);
\draw[darr] (A1b1) -- (A1b1y);
\draw[arr] (A3a2) -- (A3a2x);
\draw[darr] (B2a1) -- (B2a1x);
\draw[arr] (B2a1) -- (B2a1y);

\begin{scope}[on background layer]
  \fill[fsblue!8, rounded corners=2pt] (0.85,4.08) rectangle (3.35,4.58);
  \node[font=\tiny\itshape, text=fsblue!60, rotate=90, anchor=south] at (0.72,4.33) {share $+$r$_A$};
\end{scope}

\draw[fsblue!30, line width=2.4pt, opacity=0.5]
  (R.south) -- (A.north)
  (A.south) -- (A1.north)
  (A1.south) -- (A1b.north)
  (A1b.south) -- (A1b1.north)
  (A1b1.south) -- (A1b1x.north);
\draw[fsgreen!35, line width=2.0pt, opacity=0.45]
  (A.south) -- (A3.north)
  (A3.south) -- (A3a.north)
  (A3a.south) -- (A3a2.north)
  (A3a2.south) -- (A3a2x.north);
\draw[fsorange!30, line width=2.0pt, opacity=0.45]
  (R.south) -- (B.north)
  (B.south) -- (B2.north)
  (B2.south) -- (B2a.north)
  (B2a.south) -- (B2a1.north)
  (B2a1.south) -- (B2a1y.north);

\node[font=\tiny, text=fsblue, align=left, anchor=west] at (6.15,1.45)
  {\textbf{shared prefixes}\\[-1pt]
   A1b1x $\cap$ A3a2x $=$ trunk${+}$r$_A$\\[-1pt]
   either $\cap$ B2a1y $=$ trunk only};
\node[font=\tiny, text=fsorange!80!black, align=left, anchor=west] at (6.15,0.55)
  {\textbf{threshold abort}\\[-1pt]
   $s{<}\tau$: A2, B1, A1a, A1b1y, B2a1x,\\[-1pt]
   and C's four leaves};
\node[font=\tiny, text=fspurple, align=left, anchor=west] at (6.15,-0.40)
  {\textbf{cascade abort} (fork C, depth 2)\\[-1pt]
   CoW-3: C1a/C1b/C2a/C2b all die\\[-1pt]
   $\Rightarrow$ C1, C2, then C; Root lives (A,B)};

\node[live, minimum width=0.68cm] at (0.50,-0.78) {live};
\node[dead, minimum width=0.88cm] at (1.75,-0.78) {$s{<}\tau$};
\node[casc, minimum width=1.28cm] at (3.40,-0.78) {cascade $\Uparrow$};

% ============================================================
% (b)
% ============================================================
\begin{scope}[shift={(11.55,-0.78)}]
  \node[font=\small\bfseries, anchor=west] at (0,6.55) {(b) Many paths, staged CoW, one trunk};
  \node[font=\tiny, text=fsgray, anchor=west] at (0,6.22)
    {aligned layers $=$ shared pages (refcount $>$ 1)};

  \node[font=\tiny\sffamily\bfseries, text=fsblue] at (0.85,5.88) {A1b1x};
  \node[font=\tiny\sffamily\bfseries, text=fsgreen!50!black] at (2.55,5.88) {A3a2x};
  \node[font=\tiny\sffamily\bfseries, text=fsorange!80!black] at (4.25,5.88) {B2a1y};

  \fill[fslight, draw=fsblue, rounded corners=1.2pt] (0.10,5.28) rectangle (5.00,5.70);
  \node[font=\tiny\ttfamily] at (2.55,5.49) {shared trunk $L$ \quad ref${=}3$};

  \fill[fsblue!18, draw=fsblue!60, rounded corners=1.1pt] (0.10,4.72) rectangle (3.30,5.14);
  \node[font=\tiny\ttfamily] at (1.70,4.93) {$+$r$_A$ \ ref${=}2$};
  \fill[fsspec, draw=fsorange!70, rounded corners=1.1pt] (3.45,4.72) rectangle (5.00,5.14);
  \node[font=\tiny\ttfamily] at (4.22,4.93) {$+$r$_B$};

  \foreach \y/\la/\lb/\lc in {
    4.16/{$+$r$_{A1}$}/{$+$r$_{A3}$}/{$+$r$_{B2}$},
    3.60/{$+$r$_{A1b}$}/{$+$r$_{A3a}$}/{$+$r$_{B2a}$},
    3.04/{$+$r$_{A1b1}$}/{$+$r$_{A3a2}$}/{$+$r$_{B2a1}$},
    2.48/{$+$r$_{A1b1x}$}/{$+$r$_{A3a2x}$}/{$+$r$_{B2a1y}$}} {
    \fill[fscommit, draw=fsgreen!65!black, rounded corners=1.0pt]
      (0.10,\y) rectangle (1.60,\y+0.40);
    \node[font=\tiny\ttfamily] at (0.85,\y+0.20) {\la};
    \fill[fscommit, draw=fsgreen!65!black, rounded corners=1.0pt]
      (1.80,\y) rectangle (3.30,\y+0.40);
    \node[font=\tiny\ttfamily] at (2.55,\y+0.20) {\lb};
    \fill[fsspec, draw=fsorange!70, rounded corners=1.0pt]
      (3.45,\y) rectangle (5.00,\y+0.40);
    \node[font=\tiny\ttfamily] at (4.22,\y+0.20) {\lc};
  }

  \fill[black!8, draw=fsgray, densely dashed, rounded corners=1.0pt]
    (0.10,1.78) rectangle (1.60,2.18);
  \node[font=\tiny\ttfamily, text=fsgray] at (0.85,1.98) {$+$r y};
  \fill[black!8, draw=fsgray, densely dashed, rounded corners=1.0pt]
    (3.45,1.78) rectangle (5.00,2.18);
  \node[font=\tiny\ttfamily, text=fsgray] at (4.22,1.98) {$+$r x};
  \node[font=\tiny, text=fsorange!80!black, anchor=west] at (5.12,1.98)
    {$s{<}\tau$: free only this $+$r};

  \node[font=\tiny, text=fsgray, anchor=west, align=left] at (0.10,1.42)
    {A1b1x and A3a2x CoW five times each, but \textbf{share} trunk and $+$r$_A$.};

  \fill[fsorange!12, draw=fsorange!70, rounded corners=1.4pt]
    (0.10,0.42) rectangle (6.55,1.20);
  \node[font=\tiny, anchor=west, align=left] at (0.20,0.81)
    {\textbf{Not} 3$\times$ full clones, and not one CoW for the whole tree.\\[-1pt]
     Each fork is a CoW stage. $N$ live paths sharing a prefix of length $d$\\[-1pt]
     pay $1$ trunk $+\,d$ shared residuals $+\,$private tails---not $N$ trunks.};
\end{scope}
\end{tikzpicture}
\vspace{-0.35em}
}
\caption{\textbf{Staged CoW along many paths, threshold abort, cascade abort, shared prefixes.}
(a)~Each edge is one CoW: the child aliases the parent's pages and allocates only a residual ($+$r).
Live paths A1b1x, A3a2x, and B2a1y share the trunk; the first two also share $+$r$_A$.
Fork C has depth 2: at CoW-3 all four leaves die, so abort cascades through C1 and C2, then C; it stops at Root, which still has live children A and B.
(b)~KV stacks of the three live leaves, layers aligned: a shared layer is one physical page set with refcount ${>}1$, not $N$ copies.
Aborting A1b1y or B2a1x frees only that leaf residual; cascade-aborting C also frees $+$r$_{C1}$, $+$r$_{C2}$, and $+$r$_C$.}
\label{fig:tree}
\end{figure*}

\paragraph{The cascade as an action.}
The order above is one action per residual.
Write $\Act=\{\mathrm{HS},\mathrm{DS},\mathrm{EA},\mathrm{HP},\mathrm{PF}\}$
for hash skip, draft skip, early abort, hash partial, and prefill.
Let $m_i$ be the longest prefix of $\rho_i$ covered by full pages of the published index $H$, so $m_i\in\{0,P,2P,\ldots\}$ and $m_i\le\ell_i$.
The probe key is $\mathrm{trunk}\Vert\rho_i$, but the trunk itself is the alias written by \texttt{fork} ($\sigma_i\leftarrow(\sigma_u,|x_u|)$), so $L$ is absent from $\Work_i$ and $\Work_i\le\ell_i$ on every action.
Cases below are first-match.
\begin{equation}
\label{eq:action}
a_i=\begin{cases}
\mathrm{HS} & m_i=\ell_i,\\
\mathrm{HP} & 0<m_i<\ell_i,\\
\mathrm{DS} & m_i=0\text{ and }\rho_i\text{ is a repeated loop,}\\
\mathrm{EA} & m_i=0\text{ and the first }\alpha\ell_i\text{ already fails,}\\
\mathrm{PF} & \text{otherwise }(m_i=0).
\end{cases}
\end{equation}
Winner $i{=}0$ is forced to $\Keep$.
GPU work and connector payload are
\begin{align}
\Work_i &=
\begin{cases}
0 & a_i\in\{\mathrm{HS},\mathrm{DS}\},\\
\lceil\alpha\ell_i\rceil & a_i=\mathrm{EA},\\
\ell_i-m_i & a_i=\mathrm{HP},\\
\ell_i & a_i=\mathrm{PF},
\end{cases}
\label{eq:work}\\
\Xfer_i &=
\begin{cases}
\Work_i & i\in\Keep\ \wedge\ a_i\in\{\mathrm{HP},\mathrm{PF}\},\\
0 & \text{otherwise.}
\end{cases}
\label{eq:xfer}
\end{align}
A full hit is a retrieval: $m_i=\ell_i$ reads published KV and issues no kernel ($\Work_i=\Xfer_i=0$).
A proper prefix retrieves $m_i$ tokens and prefills the uncovered tail $\ell_i-m_i$.
A miss ($m_i=0$) has an empty probe. The draft score runs, and a residual that remains is a prefill of $\ell_i$ on the aliased trunk.
\Cref{fig:dragon} is that rule on the residual ``find the dragon.''
The dragon head is the radix root on every row.
A green prefix is a published page and is spliced onto decode; an orange tail is grown by prefill.
Speculative pages stay out of $H$ until LCP commit, so the first expand of an unpublished child has $m_i=0$ and is the $\mathrm{PF}$ row.

\begin{figure*}[t]
\centering
\resizebox{\textwidth}{!}{% One residual, seven cuts. The dragon head is the radix root.
% Green = full page already in H (spliced onto decode).
% Orange = uncovered tail, grown by prefill. Gray = withheld.
\begin{tikzpicture}[
  x=1cm, y=1cm,
  font=\sffamily,
  >=Stealth,
  line cap=round, line join=round,
  chip/.style={minimum width=0.36cm, minimum height=0.40cm, inner sep=0pt,
    font=\scriptsize\ttfamily, line width=0.45pt, rounded corners=0.6pt, anchor=center},
  hit/.style={draw=fsgreen!70!black, fill=fscommit, text=black},
  sprout/.style={draw=fsorange!85!black, fill=fsspec, densely dashed, text=fsorange!45!black},
  dead/.style={draw=fsgray, fill=black!5, densely dashed, text=fsgray},
  boxg/.style={draw=fsgreen!65!black, rounded corners=1.5pt, line width=0.5pt},
  boxo/.style={draw=fsorange!80!black, rounded corners=1.5pt, densely dashed, line width=0.55pt},
  boxd/.style={draw=fsgray, rounded corners=1.5pt, densely dashed, line width=0.4pt},
  dec/.style={draw=fsblue, fill=fslight, rounded corners=1.3pt, line width=0.5pt,
    font=\tiny\sffamily\bfseries, minimum width=1.28cm, minimum height=0.46cm, inner sep=1pt},
  decx/.style={dec, draw=fsgray, fill=black!4, text=fsgray, densely dashed},
  arr/.style={-{Stealth[length=2.3pt,width=1.7pt]}, line width=0.55pt, draw=fsblue!80},
  tick/.style={draw=fsorange!90!black, line width=0.75pt},
  tag/.style={font=\scriptsize\sffamily\bfseries, anchor=west},
  rhs/.style={font=\scriptsize\sffamily, anchor=west, align=left},
  note/.style={font=\tiny\sffamily, text=fsgray},
  dhead/.pic={
    \begin{scope}[line width=0.65pt, draw=fsblue, fill=fslight]
      % neck
      \draw[fill=fslight] (-0.04,-0.11) -- (0.14,-0.07) -- (0.14,0.11) -- (-0.04,0.15) -- cycle;
      % cranium
      \draw[fill=fslight] (0.24,0.02) ellipse (0.20 and 0.16);
      % snout, pointing at the body
      \draw[fill=fslight] (0.38,0.05) .. controls (0.58,0.13) and (0.74,0.06) .. (0.80,0.00)
        .. controls (0.74,-0.07) and (0.56,-0.09) .. (0.36,-0.05) -- cycle;
      \fill[fsblue] (0.70,0.005) circle (0.020);
      % eye
      \fill[white, draw=fsblue] (0.28,0.06) circle (0.040);
      \fill[fsblue] (0.292,0.06) circle (0.020);
      % horns
      \draw (0.14,0.15) .. controls (0.08,0.28) .. (-0.04,0.36);
      \draw (0.26,0.17) .. controls (0.24,0.30) .. (0.16,0.40);
    \end{scope}
  },
]
\path (0,-0.55) rectangle (16.20,9.15);

% ----- legend -----
\pic[scale=0.72] at (0.15,8.78) {dhead};
\node[note, anchor=west] at (1.05,8.78) {radix root};
\node[chip, hit] at (2.85,8.78) {a};
\node[note, anchor=west] at (3.12,8.78) {in $H$, spliced};
\node[chip, sprout] at (5.35,8.78) {g};
\node[note, anchor=west] at (5.62,8.78) {grow (prefill)};
\node[chip, dead] at (7.85,8.78) {x};
\node[note, anchor=west] at (8.12,8.78) {withheld};
\node[dec] at (10.55,8.78) {decode};
\node[note, anchor=west] at (11.30,8.78) {splice onto decode};
\draw[fsblue!25, line width=0.3pt] (0.05,8.42) -- (16.05,8.42);

% \spell{id}{y}{i/char/style,...}
\newcommand{\spell}[3]{%
  \foreach \i/\c/\s in {#3} {%
    \pgfmathsetmacro{\xx}{3.20+\i*0.43+(\i>=4)*0.20+(\i>=7)*0.22}%
    \node[chip,\s] (t#1\i) at (\xx,#2) {\c};%
  }%
}
\newcommand{\headat}[1]{%
  \pic at (1.18,#1) {dhead};%
  \draw[line width=0.55pt, draw=fsblue!75] (1.98,#1) -- (2.95,#1);%
}
\newcommand{\pad}[2]{($(#1.north west)+(-0.04,0.05)$) rectangle ($(#2.south east)+(0.04,-0.05)$)}
\newcommand{\cutat}[3]{% id, left index, right index
  \coordinate (k#1) at ($(t#1#2.east)!0.5!(t#1#3.west)$);%
  \draw[tick] ($(k#1)+(0,-0.26)$) -- ($(k#1)+(0,0.22)$);%
  \node[font=\tiny\sffamily, text=fsorange!80!black, anchor=south] at ($(k#1)+(0,0.40)$) {$m_i$};%
}
\newcommand{\join}[3]{% id, last index, y
  \node[dec] (D#1) at (11.35,#3) {decode};%
  \draw[arr] (t#1#2.east) -- (D#1.west);%
}
\newcommand{\act}[3]{% y, name, work
  \node[rhs] at (12.55,#1) {\textbf{#2}\\[-1pt]{\tiny #3}};%
}

% (a) full hit -------------------------------------------------
\def\ya{7.55}
\node[tag] at (0.02,\ya) {(a)};
\headat{\ya}
\spell{a}{\ya}{0/f/hit,1/i/hit,2/n/hit,3/d/hit,4/t/hit,5/h/hit,6/e/hit,7/d/hit,8/r/hit,9/a/hit,10/g/hit,11/o/hit,12/n/hit}
\draw[boxg] \pad{ta0}{ta12};
\join{a}{12}{\ya}
\act{\ya}{HS}{$\Work_i{=}0$}

% (b) front half of dragon: dra | gon ------------------------------
\def\yb{6.25}
\node[tag] at (0.02,\yb) {(b)};
\headat{\yb}
\spell{b}{\yb}{0/f/hit,1/i/hit,2/n/hit,3/d/hit,4/t/hit,5/h/hit,6/e/hit,7/d/hit,8/r/hit,9/a/hit,10/g/sprout,11/o/sprout,12/n/sprout}
\draw[boxg] \pad{tb0}{tb9};
\draw[boxo] \pad{tb10}{tb12};
\cutat{b}{9}{10}
\join{b}{12}{\yb}
\act{\yb}{HP}{$\Work_i{=}\ell_i{-}m_i$}

% (c) cut after the first letter of dragon -------------------------
\def\yc{4.95}
\node[tag] at (0.02,\yc) {(c)};
\headat{\yc}
\spell{c}{\yc}{0/f/hit,1/i/hit,2/n/hit,3/d/hit,4/t/hit,5/h/hit,6/e/hit,7/d/hit,8/r/sprout,9/a/sprout,10/g/sprout,11/o/sprout,12/n/sprout}
\draw[boxg] \pad{tc0}{tc7};
\draw[boxo] \pad{tc8}{tc12};
\cutat{c}{7}{8}
\join{c}{12}{\yc}
\act{\yc}{HP}{$\Work_i{=}\ell_i{-}m_i$}

% (d) cut before the last character --------------------------------
\def\yd{3.65}
\node[tag] at (0.02,\yd) {(d)};
\headat{\yd}
\spell{d}{\yd}{0/f/hit,1/i/hit,2/n/hit,3/d/hit,4/t/hit,5/h/hit,6/e/hit,7/d/hit,8/r/hit,9/a/hit,10/g/hit,11/o/hit,12/n/sprout}
\draw[boxg] \pad{td0}{td11};
\draw[boxo] \pad{td12}{td12};
\cutat{d}{11}{12}
\join{d}{12}{\yd}
\act{\yd}{HP}{$\Work_i{=}\ell_i{-}m_i$}

% (e) miss: the body grows from the neck ----------------------------
\def\ye{2.35}
\node[tag] at (0.02,\ye) {(e)};
\headat{\ye}
\spell{e}{\ye}{0/f/sprout,1/i/sprout,2/n/sprout,3/d/sprout,4/t/sprout,5/h/sprout,6/e/sprout,7/d/sprout,8/r/sprout,9/a/sprout,10/g/sprout,11/o/sprout,12/n/sprout}
\draw[boxo] \pad{te0}{te12};
\draw[tick] (2.55,\ye-0.26) -- (2.55,\ye+0.22);
\node[font=\tiny\sffamily, text=fsorange!80!black, anchor=north] at (2.15,\ye-0.30) {$m_i{=}0$};
\join{e}{12}{\ye}
\act{\ye}{PF}{$\Work_i{=}\ell_i$}

% (f) early abort: first three characters grow, then stop -----------
\def\yf{1.05}
\node[tag] at (0.02,\yf) {(f)};
\headat{\yf}
\spell{f}{\yf}{0/f/sprout,1/i/sprout,2/n/sprout,3/d/dead,4/t/dead,5/h/dead,6/e/dead,7/d/dead,8/r/dead,9/a/dead,10/g/dead,11/o/dead,12/n/dead}
\draw[boxo] \pad{tf0}{tf2};
\draw[boxd] \pad{tf3}{tf12};
\coordinate (kf) at ($(tf2.east)!0.5!(tf3.west)$);
\draw[tick] ($(kf)+(0,-0.26)$) -- ($(kf)+(0,0.22)$);
\node[font=\tiny\sffamily, text=fsorange!80!black, anchor=south] at ($(kf)+(0,0.40)$) {$\alpha$};
\node[decx] (Sf) at (11.35,\yf) {stop};
\draw[densely dashed, draw=fsgray, line width=0.45pt] ($(kf)+(0.06,0)$) -- (Sf.west);
\act{\yf}{EA}{$\Work_i{=}\lceil\alpha\ell_i\rceil$}

% (g) draft skip: nothing grows -------------------------------------
\def\yg{-0.25}
\node[tag] at (0.02,\yg) {(g)};
\headat{\yg}
\spell{g}{\yg}{0/f/dead,1/i/dead,2/n/dead,3/d/dead,4/t/dead,5/h/dead,6/e/dead,7/d/dead,8/r/dead,9/a/dead,10/g/dead,11/o/dead,12/n/dead}
\draw[boxd] \pad{tg0}{tg12};
\node[font=\scriptsize\sffamily\bfseries, text=fsred] at (10.55,\yg) {$\times$};
\node[note] at (10.55,\yg-0.32) {no decode};
\act{\yg}{DS}{$\Work_i{=}0$}

% word guides, drawn once over row (a) so the cut inside dragon is readable
\node[note] at (3.85,8.18) {find};
\node[note] at (5.55,8.18) {the};
\node[note] at (7.70,8.18) {dragon};
\end{tikzpicture}
}
\caption{\textbf{Seven cuts of ``find the dragon.''}
The dragon head is the radix root; the body is one residual.
Green tokens are full pages in $H$ and are spliced onto decode.
Orange tokens are absent from $H$ and are grown by prefill.
Gray tokens are withheld.
Chips stand in for tokens: a real hit ends on a full page, and each row is a different value of $m_i$.
(a)~The body is a hit ($\mathrm{HS}$, $m_i=\ell_i$, $\Work_i=0$).
(b)~The front half ``dra'' is in the tree and ``gon'' grows ($\mathrm{HP}$, $\Work_i=\ell_i-m_i$).
(c)~The cut falls after ``d''.
(d)~The cut falls before the last character.
(e)~A miss grows the body from the neck ($\mathrm{PF}$, $m_i=0$, $\Work_i=\ell_i$); the head stays the alias.
(f)~The first $\lceil\alpha\ell_i\rceil$ tokens grow and the kernel stops ($\mathrm{EA}$).
(g)~A repeated loop is withheld ($\mathrm{DS}$, $\Work_i=0$). A low score that is not a loop is not this row: it is prefilled, then decoded for a probe (\Cref{alg:app}).}
\label{fig:dragon}
\end{figure*}
Hash-local hits stay on the decode instance (the same published blocks are APC-visible there).
Draft-skipped and early-aborted residuals never enter the buffer.
\Cref{alg:app} is the per-fan-out plan; \Cref{tab:span} is the first-fan-out span.

\begin{algorithm}[t]
\caption{PrefillPlan($\{\rho_i\}$, $H$, $\tau$, $P$, winner)}
\label{alg:app}
\small
\begin{algorithmic}[1]
\REQUIRE residuals, hash index $H$, threshold $\tau$, probe length $P$
\ENSURE actions $a_i$, work $\Work_i$, transfer $\Xfer_i$
\FOR{$i=1$ \TO $k$}
  \STATE $m_i\leftarrow H.\mathrm{lookup}(\mathrm{trunk}\Vert\rho_i)$ \COMMENT{full pages of $\rho_i$}
  \IF{$m_i=\ell_i$}
    \STATE $a_i\leftarrow\mathrm{HS}$; $\Work_i\leftarrow 0$ \COMMENT{retrieve}
  \ELSIF{$0<m_i<\ell_i$}
    \STATE $a_i\leftarrow\mathrm{HP}$; $\Work_i\leftarrow\ell_i-m_i$
  \ELSE
    \STATE $s_i\leftarrow\mathrm{DraftScore}(\rho_i)$ \COMMENT{$m_i=0$, no row to read}
    \IF{$i=\mathrm{winner}$ \OR $s_i\ge\tau$}
      \STATE $a_i\leftarrow\mathrm{PF}$; $\Work_i\leftarrow\ell_i$ \COMMENT{decode to budget $D$}
    \ELSIF{$\mathrm{IsLoop}(\rho_i)$}
      \STATE $a_i\leftarrow\mathrm{DS}$; $\Work_i\leftarrow 0$ \COMMENT{repeated loop}
    \ELSE
      \STATE $a_i\leftarrow\mathrm{PF}$; $\Work_i\leftarrow\ell_i$ \COMMENT{probe $P$ tokens}
      \IF{$\mathrm{DraftScore}(y_{1:P})\ge\tau$}
        \STATE continue decode to $D$
      \ELSE
        \STATE stop at $P$
      \ENDIF
    \ENDIF
  \ENDIF
  \IF{$i=\mathrm{winner}$ \OR $a_i\in\{\mathrm{HS},\mathrm{HP},\mathrm{PF}\}$}
    \STATE keep $i$
  \ENDIF
  \STATE gate \texttt{insert} by \eqref{eq:xfer}
\ENDFOR
\end{algorithmic}
\end{algorithm}

\begin{table*}[t]
\centering
\caption{Work, transfer, and memory spans (bottleneck in \textcolor{fsred}{red}).
$W_{\mathrm{pre}}$ is GPU prefill tokens; $X$ is connector tokens; $M$ is live KV after abort;
$T_{\mathrm{ab}}$ is abort cost in pages; replay is a later \texttt{open} of the published winner.
$\mathrm{PAR}$ is the number of trunk copies touched on the first expand.}
\label{tab:span}
\scriptsize
\setlength{\tabcolsep}{3.2pt}
\setlength{\aboverulesep}{0.4pt}
\setlength{\belowrulesep}{0.4pt}
\begin{tabular}{@{}lcccccc@{}}
\rowcolor{tblhead}
\toprule
 & $W_{\mathrm{pre}}$ & $X$ & $M$ after abort & $T_{\mathrm{ab}}$ & $\mathrm{PAR}_{\mathrm{trunk}}$ & Replay $W_{\mathrm{pre}}$ \\
\midrule
Recompute & \botn{$kL+\sum_i\ell_i$} & $0$ & \botn{$kL+\sum_i\ell_i$} & $\Theta(kL/P)$ & $k$ & $kL+\ell_\star$ \\
\rowcolor{tblalt}
APC (first expand) & \botn{$kL+\sum_i\ell_i$} & $0$ & $kL+\sum_i\ell_i$ & LRU of hashed $\rho$ & $k$ & $0$ \\
hash prefill (first) & \botn{$kL+\sum_i\ell_i$} & $0$ & $kL+\sum_i\ell_i$ & LRU & $k$ & $0$ \\
\rowcolor{tblalt}
P/D disagg. & \botn{$kL+\sum_i\ell_i$} & \botn{$kL+\sum_i\ell_i$} & $kL+\sum_i\ell_i$ & full KV transfer & $k$ & $0$ \\
FS CoW & $L+\sum_i\ell_i$ & $0$ & \win{$L+\ell_\star$} & $\Theta(\ell_i/P)$ & $1$ & $L_{\mathrm{miss}}$ \\
\rowcolor{tblalt}
FS+ + P/D & \win{$L+\sum_{i\in\Keep}\ell_i$} & \win{$\sum_{i\in\Keep}\Work_i$} & \win{$L+\ell_\star$} & lazy, residual only & $1$ & \win{$0$} \\
% FS+ + stop & $L+\sum_{i\in\Keep}\ell_i$ & $\sum_{i\in\Keep}\Work_i$ & \win{$L+\ell_\star$} & lazy & $1$ & $0$ \\
\bottomrule
\end{tabular}

\vspace{0.25em}
{\scriptsize
$P{=}16$.
Known-suffix items ($p{=}1$) are admitted before residuals ($p{<}1$).
Spec pages are never hashed; LCP commit publishes full blocks of $x_\star$.
% Decode stop changes $D^{+}$, not $W_{\mathrm{pre}}$.
}
\end{table*}

\begin{lemma}[First-fan-out dominance]
\label{lem:dom}
On a cold index, $\Work_{\mathrm{APP}}\le\Work_{\mathrm{CoW}}\le\Work_{\mathrm{APC}}=\Work_{\mathrm{PD}}$ and $\Xfer_{\mathrm{APP}}\le\Xfer_{\mathrm{PD}}$, with strict inequality in $\Work$ whenever at least one repeated loop is draft-skipped, and strict inequality in $\Xfer$ whenever $\Keep\neq[k]$ or any kept residual is a hash hit.
\end{lemma}

\paragraph{Block-hash pruning.}
The first stage is a table probe and issues no model FLOPs.
The index key is $\mathrm{hash}(\mathrm{parent},\;\mathrm{block\_tokens},\;\mathrm{extra})$, with $\mathrm{extra}$ carrying LoRA id, multimodal hash, or $\mathrm{cache\_salt}$ on block $0$~\cite{vllm}.
Only full pages enter $H$, which is why \eqref{eq:action} tests $m_i$ in steps of $P$.
LCP commit inserts full blocks of $x_\star$, which is what makes a later session a hash hit ($m_i=\ell_i$) rather than another fan-out.
\texttt{fork} itself does not walk hashes: aliasing stays $O(1)$, and the lookup is a filter on the residual, not part of creating the child.
A hash-local hit is also not a transfer.
The same published blocks are already APC-visible on the decode instance, so $\Xfer_i{=}0$ even though the residual is kept.

\paragraph{Draft-score rejection.}
On a miss ($m_i=0$ in \eqref{eq:action}) the next stage is a CPU score, still before the prefill kernel.
Each residual receives a score from illegal-token and repetition checks, token entropy, or an optional draft model.
An earlier draft of this rule treated every score below $\tau$ as a draft skip.
That rule is not admissible.
A residual the text score marks illegal still solves contest items once it is decoded, while a repeated loop solves none (\Cref{tab:grow}).
Draft skip is therefore only the repeated loop: that residual is not submitted to prefill and is not inserted as a hash key.
The default prefill threshold is $\tau_{\mathrm{pre}}{=}0.15$ for ESC, Speculative Rejection, and DPTS.
The winner, and every other residual whose score is at least $\tau_{\mathrm{pre}}$, is prefilled and decoded to the budget $D$.
A non-winner that scores below $\tau$ and is not a loop is still prefilled.
It is then decoded for a probe of $P$ tokens, and the same score is applied to the generated prefix $y_{1:P}$.
If that score is at least $\tau$, decode continues from $P$ to $D$; otherwise it stops at $P$.
The probe is a prefix of decode, not the residual early-abort in panel~(f) of \Cref{fig:dragon}.
On the measurements below, $\tau{=}0.45$ and $P\in\{64,128\}$.
The control-plane rows in \S\ref{sec:eval-app} still charge the earlier blanket skip, which drops both low-score residuals before prefill.

\paragraph{Gated connector insert.}
Prefill--decode disaggregation places prefill and decode on separate instances~\cite{distserve,splitwise}.
It does not choose a branch.
The connector stays vLLM-shaped---\texttt{insert} is non-blocking, \texttt{drop\_select} blocks---and a production pipe (Nixl, LMCache, Mooncake) keeps the same gate.
Baseline disaggregation uses \eqref{eq:xfer-stock} and therefore transfers the computed span, including the $k$-way trunk.
\app uses \eqref{eq:xfer}: $\Xfer_i$ is nonzero only for a kept residual whose action is $\mathrm{HP}$ or $\mathrm{PF}$, and then only for the miss tail that the kernel computed.
Rejected residuals and hash-local hits stay on the prefill instance, so gated transfer reduces both $\Work$ and $\Xfer$ relative to baseline disaggregation.
Let $\mu$ be microseconds per transferred token.
Fan-out time is
\begin{equation}
\label{eq:tfan}
T_{\mathrm{fan}}=\Pref(\Work)+\mu\,\Xfer+T_{\mathrm{CoW}}+T_{\mathrm{abort}},
\end{equation}
with $T_{\mathrm{CoW}}$ a pointer swap (no KV memcpy) and $T_{\mathrm{abort}}$ deferrable when \texttt{abort} is lazy.
Under the control-plane model of \S\ref{sec:eval-app} ($\ell{\times}12\,\mu\mathrm{s}$, $\mu{=}2\,\mu\mathrm{s}$), \eqref{eq:tfan} is $139$\,ms for APC, $162$\,ms for baseline disaggregation, $47$\,ms for CoW, and $31.4$\,ms for \app.

\iffalse
\paragraph{Orthogonal decode-length stop.}
The cascade changes $\Work$ and $M$.
End-to-end time is still $T_{\mathrm{fan}}+T_{\mathrm{dec}}(D)$, so a shorter prefill does not by itself shorten decode.
ForkServe+ stops committed decode only when the last non-empty line is a finished answer (\texttt{\#\#\#\#} or a closed \texttt{\textbackslash boxed}).
An earlier \texttt{\#\#\#\#} is an intermediate step, and a bare \texttt{</think>} is not an answer.
If item $j$ hits a marker at $t_j$, emitted tokens are $D^{+}=\sum_j\min(t_j,D)\le nD$, and $T_{\mathrm{dec}}$ scales with $D^{+}$ when decode is memory-bandwidth bound.
The stop ends a finished answer; it does not draft output tokens, and it does not change which residual was admitted.
\fi

\subsection{The Spine on the Critical Path}
\label{sec:schedule}
After \app, fan-out is the prefill of the kept residual.
End-to-end time is still $T_{\mathrm{fan}}+T_{\mathrm{dec}}$, and $T_{\mathrm{dec}}$ is the long term.
Three placements put the spine on that sum.
A fixed chunk width and a fan-out barrier leave the freed KV and the pruned queue off the path.

\paragraph{Slots follow the spine.}
Let $H$ be the KV pool after weights.
A live tree occupies $M_{\mathrm{spine}}=b(L+\ell_\star)$.
APC still holds the aborted residuals until they expire, so its live set is the larger $M_{\mathrm{APC}}$.
The number of concurrent decode slots is
\begin{equation}
\label{eq:slots}
C=\bigl\lfloor H/M\bigr\rfloor,
\end{equation}
with $M=M_{\mathrm{spine}}$ for \sys and $M=M_{\mathrm{APC}}$ for the prefix cache.
$C$ is not a constant chunk size.
A request that finishes its current decode step releases the slot, and the next queued spine takes it.
Throughput is
\begin{equation}
\label{eq:theta}
\Theta\;\propto\; C/T_{\mathrm{e2e}}.
\end{equation}
The same $H$ therefore admits more sessions once occupancy is the spine, and the wall clock of a batch is the number of waves $n/C$, not $n$ divided by a fixed width.

\paragraph{Prefill the next winner during decode.}
Each tick $t$ has a token budget $B_t$ from the chunked-prefill cap~\cite{sarathi}.
JITServe's margin~\cite{jitserve} gives each committed request $r$ a minimum token rate $\underline{s}_r$.
\begin{align}
B^{\mathrm{c}}_t &= \min\Bigl(B_t,\; \sum_{r\in R^{\mathrm{c}}_t}\max(\underline{s}_r\Delta,\;1)\Bigr), \label{eq:bc}\\
B^{\mathrm{s}}_t &= B_t-B^{\mathrm{c}}_t. \label{eq:bs}
\end{align}
$B^{\mathrm{c}}_t$ is committed decode.
$B^{\mathrm{s}}_t$ is not spent on the $k{-}1$ residuals \app already rejected.
It prefills the miss tail of the next winner, in the same step as the running decode.
If $B^{\mathrm{s}}_t=0$, no speculative chunk is admitted.
The fan-out of losers is not a generate that must finish before any decode token.
A tree-sticky router pins the trunk to the worker that prefills the root; only residual work may migrate.

\paragraph{Hide the known suffix in the tool pause.}
A candidate with $p{=}1$ is an exact wrapper, not a guess.
If $\Pref(|w|)\le\Idle(u)$, that prefill runs while the tool is in flight, including when the parent is still marked committed.
The pause is the deadline.
A suffix that would finish after $\Idle(u)$ is rejected, as is any observation residual the pause cannot cover.
When the observation arrives, LCP commit checks the pinned prefix and the committed prefill is the unmatched tail.
Under P/D the same suffix is prefiller work tagged $\texttt{soft-deadline}=\Idle(u)$ and is inserted once, onto the trunk owner.
Speculative tokens still do not enter the sampler, so Theorem~\ref{thm:correct} is unchanged.
